This report asked whether a patch-based forecaster preserves the frequency content of signals that classical sampling theory says are fully representable, and how $f_s$, $P$ and $S$ decide when it does not. The characterisation is complete and depends on no fit. A tone completing a whole number of cycles in one stride repeats its tokens; one completing a whole number in one patch repeats their content; and the union of the two combs is a closed-form candidate set computable for a geometry never measured. That set locates where a loss could sit, not that one exists, and the report holds the two apart throughout.

Whether a trained model does lose information there remains open, and the design that decides it is fixed in advance: it returns not reportable or not identified rather than a null when a fit fails to earn a verdict. Two consequences already follow: patch geometry is a deployment parameter to be chosen against the known spectral content of a plant, not a default to inherit; and since a detector placed after tokenisation would inherit any blind spot there, monitoring must act on the raw signal upstream. Neither waits on the fits, and neither is available to a system that treats the tokeniser as opaque.

The exploratory steps agree: added tones are recovered more often on the sites than beside them, and on the photovoltaic series no retrained geometry beats the $P=S=16$ baseline, with no appreciable energy at the sites. The Bayesian analysis returned a direction but no verdict. Before the gates, the rules would refute $H1$ on both readings: at a candidate frequency the forecast recovers the tone $35\%$ better than beside it, and the internal state needs only $6\%$ more code to read it. They would support $H2$, the effect barely changing with phase, and $H3a$, the stride sites moving exactly as $f_s/S$ predicts; the patch-grid dip lies inside its equivalence band, so $H3$ location would be refuted, $M1$ would be inconclusive and $H3b$ is not identified. Every model converged and none depends on its prior, but none reproduces the dispersion of the data and Models~A to~C and the patch branch of D2 also fail parameter recovery, so under the rules fixed in advance every claim is not reportable. For the security reading this removes no exposure and establishes none: the sites are real, predictable and treated differently by the model, which is what an operator needs to monitor them, and the direction of the evidence is that they attract the forecast rather than blind it. Until a claim is reportable, the agent of \autoref{sec:agent} keeps every candidate frequency \valueof{PartiallyObservable} and acts in \valueof{ConservativeMode}.

If confirmed, structural aliasing in a patch-based forecaster would be a configuration-determined, attacker-independent weakness whose candidate sites are derivable in closed form rather than conjectured. What the arithmetic settles on its own is that most sites exclusive to one branch lie on the patch branch (\autoref{tab:branchSites}), that the pre-registered mitigation cannot reach that branch whatever the fit returns, and that any exposure there cannot be managed by overlap inside the model. It has to be managed around the model, with the candidate set recomputed whenever the configuration changes.

Four extensions would sharpen the claim. The first is the one every verdict now waits on: models that reproduce the dispersion of the contrast, with a residual scale per geometry and a skewed likelihood. The others are independently trained instances per geometry, which a single seed cannot separate from geometry, an untrained-network baseline on the same dispersion measurement, and a test of whether a candidate set derived at $512$\,Hz survives translation to telemetry sampled once a minute.
